\section{Auto Research：自动化研究的起步}

\subsection{核心理念：移除人类瓶颈}

Karpathy 的核心主张：\textbf{要最大化利用当前工具，必须把自己从循环中移除}。不能在那里等着 prompt 下一步，必须让系统完全自主运行。目标是最大化 token 吞吐量并将人类移出 loop。

\begin{importantbox}{Auto Research 的定义}
给定一个目标、一个评估指标和操作边界，让 agent 自动循环探索改进方案。人类只需"arrange it once and hit go"。核心问题是：如何让更多 agent 运行更长时间、更少需要人类参与。
\end{importantbox}

\subsection{DataChat GPT-2 实验}

Karpathy 用他的 DataChat 项目（GPT-2 训练 playground）进行了 auto research 实验。他已经以传统方式（手动超参搜索、两个十年的经验）将模型调优到自认为不错的状态，然后让 auto researcher 运行一晚：

\begin{itemize}
    \item Auto researcher 发现了他遗漏的 weight decay on value embeddings
    \item Adam betas 未充分调优
    \item 这些超参之间存在联合交互——调一个可能需要重新调其他的
\end{itemize}

\begin{warningbox}{人类研究者的过度自信}
Karpathy 指出，拥有两个十年训练经验的研究者也会犯遗漏错误。更关键的是，人类不应该成为超参搜索的瓶颈——"I shouldn't be looking at the results. There's objective criteria."
\end{warningbox}

\subsection{Program.md：元优化}

Auto research 的配置通过一个 \texttt{program.md} 文件描述——"研究组织的 DNA"。不同的 program.md 会产生不同的研究进展。

\begin{importantbox}{研究组织即代码}
"A research organization is a set of markdown files that describe all the roles and how the whole thing connects." 可以想象对 program.md 本身进行优化——100\% 可以观察改进来自哪里，然后调整 program.md 使更多类似改进被发现。这是 meta-optimization 的层级。
\end{importantbox}

\subsection{Auto Research 的适用边界}

\begin{warningbox}{Auto Research 的两大限制}
\begin{enumerate}
    \item \textbf{需要可验证的指标}：写 CUDA kernel（同样行为但更快）是完美场景；但"用户意图理解""产品品味"等软性目标无法 auto research——"if you can't evaluate, you can't auto research it"
    \item \textbf{模型仍然 rough around the edges}：如果走得太远太快，系统反而不再有用。Agent 在一些明显问题上浪费大量 compute，让人沮丧
\end{enumerate}
\end{warningbox}

\subsection{本章小结}

Auto research 代表了从"人类在循环中手动研究"到"人类配置一次、agent 自动循环"的范式转换。它在可验证指标的领域已经展示出令人惊喜的效果，但对软性目标和模型 jaggedness 的挑战仍需正视。


